\labsection{5}{Neobank unit economics}{sec:lab05-unit-economics}

\begin{labproject}{Find out whether a digital-only bank makes money per customer}
\textbf{Coverage.} Chapter~\ref{ch:digital-banks}.\\
\textbf{Goal.} Build the model that answers the chapter's third strategic priority --- expand \emph{profitably} --- and find the assumption the whole business rests on.

\begin{labmode}
Chapter~\ref{ch:digital-banks} observes that neobanks grow customers quickly and reach profitability slowly. This lab shows why those two facts are connected rather than contradictory.

Every input below is a parameter you will vary. The point is not to produce one answer; it is to find which input the answer is most sensitive to.
\end{labmode}

\medskip\noindent\textbf{Part A --- The four quantities that decide everything.}\par

\begin{mltable}
\begin{tabularx}{\linewidth}{@{}>{\bfseries\leavevmode\color{MLNavy}}p{0.22\linewidth}X p{0.20\linewidth}@{}}
\toprule
\rowcolor{MLPaleBlue!92}
Quantity & Definition & Typical neobank \\
\midrule
CAC & Total acquisition spend $\div$ customers acquired & \pounds20--\pounds60 \\
ARPU & Revenue per user per month & \pounds2--\pounds8 \\
Cost to serve & Variable cost per user per month & \pounds1--\pounds3 \\
Churn & Monthly share of customers who leave & 2--5\% \\
\bottomrule
\end{tabularx}
\end{mltable}

From these, two derived numbers decide whether the business exists:

\[
\text{Contribution} = \text{ARPU} - \text{cost to serve}, \qquad
\text{Payback (months)} = \frac{\text{CAC}}{\text{Contribution}}
\]

\[
\text{LTV} = \frac{\text{Contribution}}{\text{monthly churn}}
\]

\begin{workedexample}
CAC \pounds45, ARPU \pounds4.20, cost to serve \pounds1.80, churn 3\%/month.

\smallskip
Contribution $= 4.20 - 1.80 = \pounds2.40$ per month.\\
Payback $= 45 \div 2.40 = 18.75$ months.\\
LTV $= 2.40 \div 0.03 = \pounds80$.\\
LTV/CAC $= 80 \div 45 = 1.78$.

\smallskip
The customer is profitable, but only just. A ratio of 1.78 leaves almost nothing for fixed costs --- engineering, compliance, licensing --- which this model has not yet touched. The usual venture benchmark is 3.0.
\end{workedexample}

\begin{lstlisting}[style=mlpython]
def unit_economics(cac, arpu, cost_to_serve, monthly_churn):
    contribution = arpu - cost_to_serve
    if contribution <= 0:
        return {"verdict": "loses money on every customer"}

    return {
        "contribution": round(contribution, 2),
        "payback_months": round(cac / contribution, 1),
        "ltv": round(contribution / monthly_churn, 2),
        "ltv_cac": round((contribution / monthly_churn) / cac, 2),
        "avg_lifetime_months": round(1 / monthly_churn, 1),
    }

print(unit_economics(cac=45, arpu=4.20, cost_to_serve=1.80, monthly_churn=0.03))
\end{lstlisting}

\begin{pitfall}
$\text{LTV} = \text{contribution} / \text{churn}$ assumes churn is constant forever, which it is not --- it is high early and falls as customers settle.

The formula therefore \emph{understates} the value of a customer who survives the first year and \emph{overstates} the value of a book dominated by new sign-ups. A neobank growing quickly has a young book, so its reported LTV is the least reliable number in its deck. Say so when you quote it.
\end{pitfall}

\medskip\noindent\textbf{Part B --- Growth can destroy a profitable business.}\par
Each new customer costs CAC today and repays it over the payback period. So acquisition consumes cash \emph{now} against profit \emph{later}, and the faster you grow, the deeper the hole:

\begin{lstlisting}[style=mlpython]
def project(months, new_per_month, cac, arpu, cost_to_serve, churn):
    customers, cash = 0, 0.0
    for m in range(1, months + 1):
        customers = customers * (1 - churn) + new_per_month
        cash += customers * (arpu - cost_to_serve) - new_per_month * cac
        if m % 6 == 0:
            print(f"month {m:3d}  customers {customers:9,.0f}  cumulative cash {cash:12,.0f}")
    return cash

project(36, new_per_month=50_000, cac=45, arpu=4.20, cost_to_serve=1.80, churn=0.03)
\end{lstlisting}

Run it at 50{,}000 new customers per month, then again at 5{,}000, and compare:

\begin{lstlisting}[style=mlterminal]
month 60, fast (50,000/mo):  1,398,656 customers    -3,535,671 cash
month 60, slow ( 5,000/mo):    139,866 customers      -353,567 cash

fast/slow cash ratio: 10.0x   break-even: neither, within 60 months
\end{lstlisting}

\begin{pitfall}
It is tempting to conclude that the slower bank is closer to profitability. It is not. The two curves are \emph{exact} multiples of each other --- 10.0$\times$, matching the 10$\times$ difference in acquisition --- and they cross zero in the same month.

Both customers and cash are linear in \texttt{new\_per\_month}, so growth rate changes the \emph{size} of the hole and not its \emph{shape}. Slowing down does not bring profitability forward; it only makes the loss smaller in absolute terms while producing proportionally fewer customers.
\end{pitfall}

So what does move break-even? Either better unit economics, or stopping acquisition altogether. Run the projection once more with \texttt{stop\_acquiring\_at=24}:

\begin{lstlisting}[style=mlterminal]
month 24:   814,305 customers   -22,940,039   <- acquisition stops
month 39:                                 0   <- cash break-even
month 60:   272,000 customers   +19,142,779
\end{lstlisting}

The book keeps paying while the CAC stops, and the business turns cash-positive at month 39 --- shedding roughly two-thirds of its customers to churn along the way.

\begin{keyidea}
This is why a neobank can be simultaneously profitable per customer and heavily loss-making overall, and why growth is funded by equity rather than earnings.

It also explains the strategic sequence the chapter describes: acquire on a free product to drive CAC down, then raise ARPU by cross-selling credit and FX. Neither move works alone. Acquisition without monetisation never repays; monetisation without scale never covers fixed costs.

And it explains what a board does when funding tightens --- cut acquisition and harvest the existing book, trading future scale for present solvency.
\end{keyidea}

\medskip\noindent\textbf{Part C --- Sensitivity: which input actually matters?}\par
Vary each of the four inputs by $\pm 20\%$ and record the effect on LTV/CAC:

\begin{lstlisting}[style=mlpython]
base = dict(cac=45, arpu=4.20, cost_to_serve=1.80, monthly_churn=0.03)
baseline = unit_economics(**base)["ltv_cac"]

for key in base:
    for factor in (0.8, 1.2):
        trial = dict(base)
        trial[key] = base[key] * factor
        result = unit_economics(**trial)["ltv_cac"]
        print(f"{key:15s} x{factor}  LTV/CAC {result:5.2f}  "
              f"({(result / baseline - 1):+.0%})")
\end{lstlisting}

\begin{lstlisting}[style=mlterminal]
baseline LTV/CAC: 1.78

input                  -20%       +20%      swing
arpu                   1.16       2.40       1.24
monthly_churn          2.22       1.48       0.74
cac                    2.22       1.48       0.74
cost_to_serve          2.04       1.51       0.53
\end{lstlisting}

ARPU dominates, with nearly double the swing of anything else --- and the reason is structural rather than accidental. Churn and CAC each enter LTV/CAC once, so a 20\% change moves the ratio by roughly 20\%. ARPU enters through \emph{contribution}, which is a difference: a 20\% rise in ARPU of \pounds4.20 is \pounds0.84, and against a contribution of only \pounds2.40 that is a 35\% increase. Subtracting a fixed cost amplifies every proportional change in revenue.

\begin{tryit}
Confirm that reasoning by setting cost to serve to zero and rerunning. The result:

\begin{lstlisting}[style=mlterminal]
cac             1.30
monthly_churn   1.30
arpu            1.24
\end{lstlisting}

Read it carefully, because the mechanism is not what it first looks like. ARPU's swing does not fall --- it stays at 1.24. The other two \emph{rise} to meet it, because removing the subtraction lifts the baseline ratio from 1.78 to 3.11 and every proportional change is now measured against a larger number.

With nothing subtracted, $\text{LTV}/\text{CAC} = \text{ARPU} / (\text{churn} \times \text{CAC})$ and all three inputs enter multiplicatively, so all three matter roughly equally. CAC and churn edge ahead only because they sit in a denominator, and $1/x$ is convex --- a 20\% cut helps slightly more than a 20\% rise hurts.

ARPU's dominance in the real model is therefore created entirely by the cost-to-serve subtraction. That is a property of the business, not of the arithmetic.

Then connect it to strategy: if that input is the one that matters, which of Chapter~\ref{ch:digital-banks}'s three strategic priorities should the bank fund first? Does your answer match what the chapter recommends?
\end{tryit}
\end{labproject}

\begin{verification}
\begin{itemize}
  \item Contribution margin is computed before LTV; a non-positive contribution is reported as such rather than producing a meaningless LTV.
  \item The constant-churn assumption behind LTV is stated whenever LTV is quoted.
  \item The cash projection is run at two growth rates and the difference explained.
  \item The sensitivity analysis names one dominant input with a reason drawn from the formula.
\end{itemize}
\end{verification}

\begin{labdeliverable}
Submit the model output, the two cash projections, the sensitivity table, and a two-page memo advising a neobank board: current LTV/CAC, months to cash break-even at present growth, the single input most worth improving, and one action you would take to improve it.
\end{labdeliverable}
